\chapter{Negative-Inversion Zero-Set No-Go}
\label{chap:negative-inversion-zero-set-no-go}

\status{Exact no-go theorem proved. The canonical negative inversion cannot be a global symmetry of the xi zero set; only finitely many xi zeros can be paired with another xi zero by that map. RH remains open.}

Chapter~\ref{chap:euler-riemann-negative-inversion} proved the canonical
negative inversion in the projective coordinate
\[
 N_u(u)=-\frac1u,
\]
and Chapter~\ref{chap:pauli-spinor-lift} established its concrete Pauli/SU(2)
lift.  The present chapter supplies the complementary spectral firewall: exact
operator symmetry does not imply invariance of the zero set of the Riemann xi
function.

\section{The induced map on the $s$-plane}

With
\[
 u=\Omega(s)=\frac{s}{1-s},
 \qquad
 s=\frac{u}{1+u},
\]
conjugating $N_u(u)=-1/u$ back to the $s$-plane gives
\begin{equation}
 \boxed{
 N_s(s)=\frac{s-1}{2s-1}.
 }
 \label{eq:g014-Ns}
\end{equation}
It is a Möbius involution on the Riemann sphere.  In the affine chart its pole
is $s=1/2$, and projectively it exchanges $1/2$ with infinity.

Subtracting the half gives the exact identity
\begin{equation}
 \boxed{
 N_s(s)-\frac12=-\frac1{4s-2}.
 }
 \label{eq:g014-half-defect}
\end{equation}
Therefore
\begin{equation}
 |s|\longrightarrow\infty
 \quad\Longrightarrow\quad
 N_s(s)\longrightarrow\frac12.
 \label{eq:g014-contraction}
\end{equation}

On the critical line,
\begin{equation}
 \boxed{
 N_s\!\left(\frac12+it\right)
 =\frac12+\frac{i}{4t},
 \qquad t\ne0.
 }
 \label{eq:g014-critical-height}
\end{equation}
Thus the negative inversion preserves the critical line as a geometric locus
while exchanging large height with reciprocal small height.

\section{The half is not a xi zero}

The positive-kernel theorem of Chapter~\ref{chap:positive-riemann-kernel}
gives
\begin{equation}
 \xi\!\left(\frac12\right)
 =\int_0^\infty\Phi(y)\,dy>0.
 \label{eq:g014-xi-half-positive}
\end{equation}
Hence
\[
 \boxed{\xi(1/2)\ne0.}
\]
In particular, the affine pole of $N_s$ is not itself a xi zero.

\section{The paired-zero set}

Let
\[
 Z_\xi=\{\rho\in\C:\xi(\rho)=0\}
\]
and define
\begin{equation}
 P_N
 =\{\rho\in Z_\xi:\xi(N_s(\rho))=0\}.
 \label{eq:g014-paired-set}
\end{equation}
The completed xi function is a non-zero entire function, so its zeros are
isolated.  Its zero set is infinite; for example, Hardy's classical theorem
gives infinitely many zeros on the critical line \citep{hardy1914}, and the
standard analytic theory is summarized in \citep{titchmarsh1986}.

\begin{theorem}[Negative-inversion paired-zero no-go]
\label{thm:g014-paired-finite}
The set $P_N$ in \cref{eq:g014-paired-set} is finite.
\end{theorem}

\begin{proof}
Assume that $P_N$ is infinite.  Since the zeros of the non-zero entire
function $\xi$ are isolated, an infinite subset of $Z_\xi$ cannot remain in a
compact subset of the plane.  Hence there exist distinct
\[
 \rho_n\in P_N,
 \qquad
 |\rho_n|\to\infty.
\]
By definition of $P_N$,
\[
 \xi(N_s(\rho_n))=0
\]
for every $n$.  But \cref{eq:g014-half-defect} gives
\[
 N_s(\rho_n)
 =\frac12-\frac1{4\rho_n-2}
 \longrightarrow\frac12.
\]
Continuity of the entire function $\xi$ therefore forces
\[
 \xi\!\left(\frac12\right)
 =\lim_{n\to\infty}\xi(N_s(\rho_n))
 =0,
\]
contradicting \cref{eq:g014-xi-half-positive}.  Thus $P_N$ is finite.
\end{proof}

\begin{corollary}[No global zero-set invariance]
\label{cor:g014-no-global-invariance}
The complete xi zero set is not invariant under $N_s$:
\[
 \boxed{N_s(Z_\xi)\ne Z_\xi.}
\]
\end{corollary}

\begin{proof}
If $N_s(Z_\xi)=Z_\xi$, then every xi zero would belong to $P_N$.  Since
$Z_\xi$ is infinite, this contradicts \cref{thm:g014-paired-finite}.
\end{proof}

\begin{corollary}[All but finitely many zeros fail to pair]
\label{cor:g014-cofinite-failure}
All but finitely many xi zeros are mapped by $N_s$ to points at which xi is
non-zero.
\end{corollary}

This conclusion is stronger than merely exhibiting a counterexample to global
invariance.  It says that negative-inversion pairing, if it occurs at all for
xi zeros, is necessarily an exceptional finite phenomenon.

\section{What survives from the operator geometry}

The no-go theorem does not weaken G012 or G013.  The identities
\[
 N_u(u)=-\frac1u,
 \qquad
 \{I,R,E,N\}\cong V_4,
\]
and the spinor lift
\[
 1\longrightarrow\{\pm I\}
 \longrightarrow Q_8
 \longrightarrow V_4
 \longrightarrow1
\]
remain exact.  What fails is only the stronger spectral assertion
\[
 \text{operator symmetry}\Longrightarrow\text{xi-zero-set symmetry}.
\]
The distinction is essential: the projective/Bloch algebra is a coordinate
and spinor structure, whereas the xi zero set is an analytic spectrum with
additional constraints.

\section{Numerical regression is not the proof}

The accompanying SOH-G014 test and receipt evaluate the first few
non-trivial zeros and verify numerically that their negative-inversion images
lie near $1/2$ and are far from xi zeros.  Those checks protect the
implementation against coordinate-sign mistakes.  The theorem itself is the
analytic contradiction in \cref{thm:g014-paired-finite}; no finite zero table
is used to establish it.

\section{Proof firewall}

Proved in SOH-G014:
\begin{itemize}
 \item the exact $s$-plane map \cref{eq:g014-Ns};
 \item the exact half-defect \cref{eq:g014-half-defect};
 \item the critical-line height transformation \cref{eq:g014-critical-height};
 \item finiteness of the paired-zero set $P_N$;
 \item failure of global xi-zero-set invariance;
 \item cofinite failure of negative-inversion zero pairing.
\end{itemize}

Not proved or claimed in SOH-G014:
\begin{itemize}
 \item that $P_N$ is empty;
 \item a new functional equation for $\xi$;
 \item invariance under the Euler half-turn;
 \item PF$_\infty$;
 \item SOH-G003 real-rootedness;
 \item the Riemann Hypothesis.
\end{itemize}
